% ============================================================================
% Fichier  : partie4/P04_C12_forensique_memoire.tex : FORN (LPI342)
% Sources : NIST SP 800-86 ; RFC 3227. Opération TSANAGA.
% ============================================================================

\chapter{Forensique Mémoire}
\label{chap:memoire}
\chaptermark{Forensique mémoire}

\epigraph{%
    « La mémoire vive avoue ce que le disque dissimule. »%
}{Principe du cours}

\niveauChapitre{Expert}{Chapitres~\ref{chap:acquisition} et~\ref{chap:disque}.}

\begin{boxobjectifs}
\begin{itemize}
    \item Comprendre ce que révèle la mémoire vive.
    \item Analyser une capture mémoire.
    \item Repérer processus, connexions, secrets.
    \item Relier la mémoire au disque et au réseau.
    \item Qualifier un élément de mémoire.
\end{itemize}
\end{boxobjectifs}

\bigskip

% ============================================================================
\section*{Question directrice : que contenait la mémoire de l'ordinateur de M. NDONGO au moment de la saisie ?}
% ============================================================================

La capture mémoire réalisée à chaud (Chapitre~\ref{chap:acquisition}) contient
l'état vivant du système~: c'est là que se trouvait la clé du conteneur
chiffré. Ce chapitre apprend à l'analyser.

\medskip
\textbf{Plan du chapitre.}
\begin{enumerate}
    \item Ce que révèle la mémoire (section~\ref{sec:c12-revele}).
    \item Analyser une capture (section~\ref{sec:c12-analyser}).
    \item Secrets et chiffrement (section~\ref{sec:c12-secrets}).
    \item Recouper avec le disque (section~\ref{sec:c12-recouper}).
\end{enumerate}

\separateur

% ============================================================================
\section{Ce que Révèle la Mémoire}
\label{sec:c12-revele}
% ============================================================================

\begin{tableaucours}{L{4.0cm}L{8.2cm}}{%
    \tableauHeader{Élément} &
    \tableauHeader{Intérêt}
}
Processus en cours & Programmes actifs, parent suspect \\
Connexions réseau & Destinations contactées au moment de la saisie \\
Secrets déchiffrés & Clés, mots de passe, jetons en clair \\
Fichiers ouverts & Documents en cours d'usage \\
Fragments & Textes, adresses, en mémoire mais non sur disque \\
\end{tableaucours}
\vspace{-0.8em}
\begin{center}{\small\itshape Tableau~12.1. Ce que révèle la mémoire vive}\end{center}

% ============================================================================
\section{Analyser une Capture}
\label{sec:c12-analyser}
% ============================================================================

\begin{boxprocedure}
\textbf{Analyser une capture mémoire}
\begin{enumerate}
    \item Vérifier l'intégrité de la capture (condensat).
    \item Lister les processus et leurs liens de parenté.
    \item Lister les connexions réseau ouvertes.
    \item Rechercher secrets, clés, fragments pertinents.
    \item Recouper avec le disque et le réseau.
\end{enumerate}
\end{boxprocedure}

% ============================================================================
\section{Secrets et Chiffrement}
\label{sec:c12-secrets}
% ============================================================================

\begin{boxscenario}{Opération TSANAGA. La clé en mémoire}
\textbf{Constatation.} La capture mémoire contient la clé du conteneur
chiffré, en clair (le conteneur était monté à la saisie). Elle permet d'ouvrir
le conteneur et d'en analyser le contenu.

\medskip
\textbf{Analyse} (confiance élevée). Sans la capture à chaud, le conteneur
serait resté inaccessible. La clé récupérée a permis d'établir que le
conteneur contenait une copie du fichier clients. L'arbitrage de l'acquisition
à chaud trouve ici sa justification.
\end{boxscenario}

\begin{boxCRO}
\textbf{Une clé en mémoire, un pouvoir à encadrer.} Déchiffrer un conteneur
grâce à une clé trouvée en mémoire reste dans le périmètre de la réquisition~:
\cro{0}{+}{+}. En sortir (déchiffrer d'autres supports non visés) violerait la
confidentialité et le mandat.
\end{boxCRO}

% ============================================================================
\section{Recouper avec le Disque}
\label{sec:c12-recouper}
% ============================================================================

La mémoire, seule, est un instantané~; recoupée avec le disque (le fichier
récupéré) et le réseau (les destinations), elle devient un faisceau. La force
de l'enquête tient à la convergence des domaines sur une même histoire.

\begin{boxmetier}{Analyste forensique}
\textbf{La mémoire, fenêtre sur l'instant.} L'analyste traite la capture
mémoire comme une photographie de l'instant de la saisie~: précieuse, unique,
à recouper. Il en extrait ce qui sert l'affaire, dans le périmètre.
\end{boxmetier}

\separateur

\begin{boxresume}
\textbf{Ce qu'il faut retenir de ce chapitre~:}
\begin{itemize}
    \item \textbf{La mémoire} révèle processus, connexions, secrets,
          fragments.
    \item \textbf{On vérifie} l'intégrité de la capture avant de l'analyser.
    \item \textbf{Les secrets en clair} (clés) ouvrent les conteneurs
          chiffrés.
    \item \textbf{On reste dans le périmètre} de la réquisition.
    \item \textbf{On recoupe} mémoire, disque, réseau.
\end{itemize}
\end{boxresume}

\begin{boxinfo}
\textbf{Points de contrôle de l'élève.} Je sais~: citer ce que révèle la
mémoire~; analyser une capture~; exploiter un secret dans le périmètre~;
recouper les domaines~; qualifier un élément de mémoire.
\end{boxinfo}

\bigskip

\begin{boxexo}[title={\ding{45}\quad Exercice~12.1. Révélations \niveaubase}]
Citez trois éléments que révèle la mémoire vive.
\end{boxexo}

\begin{boxexo}[title={\ding{45}\quad Exercice~12.2. Intégrité \niveaubase}]
Que vérifie-t-on avant d'analyser une capture mémoire~?
\end{boxexo}

\begin{boxexo}[title={\ding{45}\quad Exercice~12.3. La clé \niveaubase}]
Qu'a permis la clé trouvée en mémoire dans TSANAGA~?
\end{boxexo}

\begin{boxexo}[title={\ding{45}\quad Exercice~12.4. Périmètre \niveaumoyen}]
Peut-on utiliser une clé trouvée en mémoire pour déchiffrer un support non
visé~? Pourquoi~?
\end{boxexo}

\begin{boxexo}[title={\ding{45}\quad Exercice~12.5. Instantané \niveaumoyen}]
Pourquoi la mémoire est-elle un instantané à recouper~?
\end{boxexo}

\begin{boxexo}[title={\ding{45}\quad Exercice~12.6. CRO \niveaumoyen}]
Expliquez l'analyse CRO de l'usage d'une clé trouvée en mémoire.
\end{boxexo}

\begin{boxexo}[title={\ding{45}\quad Exercice~12.7. Parenté \niveaumoyen}]
Pourquoi examiner les liens de parenté entre processus~?
\end{boxexo}

\begin{boxexo}[title={\ding{45}\quad Exercice~12.8. Convergence \niveaumoyen}]
Comment la mémoire se recoupe-t-elle avec le disque dans TSANAGA~?
\end{boxexo}

\begin{boxexo}[title={\ding{45}\quad Exercice~12.9. Plan d'analyse \niveauexpert}]
Rédigez le plan d'analyse de la capture mémoire de TSANAGA.
\end{boxexo}

\begin{boxexo}[title={\ding{45}\quad Exercice~12.10. Justification \niveauexpert}]
Rédigez la note qui justifie l'exploitation de la clé trouvée en mémoire.
\end{boxexo}

\bigskip

\begin{boxinfo}
\textbf{Pour aller plus loin.}
\begin{itemize}
    \item NIST SP~800-86~\cite{nist80086}~; RFC~3227~\cite{rfc3227}.
\end{itemize}
\end{boxinfo}

\begin{boxinfo}
\textbf{Transition.} Reste la trace des échanges~: le
Chapitre~\ref{chap:reseau} traite du forensique réseau.
\end{boxinfo}
